\documentclass[11pt,a4paper]{article}
\usepackage[T1]{fontenc}
\usepackage[utf8]{inputenc}
\usepackage{amsmath,amsthm,mathtools}
\usepackage{newpxtext,newpxmath}
\usepackage[margin=27mm,headheight=15pt]{geometry}
\usepackage{microtype}
\usepackage{booktabs,array,longtable}
\usepackage{xcolor}
\usepackage{enumitem}
\usepackage{fancyhdr,titlesec}
\usepackage[most]{tcolorbox}
\usepackage{listings}
\usepackage{hyperref}
\usepackage[nameinlink,noabbrev]{cleveref}
\definecolor{oliveink}{HTML}{4B573A}
\definecolor{olivewash}{HTML}{F3F4ED}
\definecolor{mutedink}{HTML}{62665D}
\definecolor{ruleink}{HTML}{B8BDA8}
\hypersetup{colorlinks=true,linkcolor=oliveink,citecolor=oliveink,urlcolor=oliveink,
  pdftitle={When the Secant Numbers Become Periodic},
  pdfauthor={Research report prepared for Vladimir Reshetnikov},
  pdfsubject={A correction to an OEIS conjecture, exact periods and preperiod densities}}
\urlstyle{same}
\titleformat{\section}{\Large\bfseries\color{oliveink}}{\thesection}{0.65em}{}
\titleformat{\subsection}{\large\bfseries\color{oliveink}}{\thesubsection}{0.65em}{}
\titlespacing*{\section}{0pt}{2.5ex plus 1ex minus .2ex}{1.2ex plus .2ex}
\titlespacing*{\subsection}{0pt}{2ex plus .7ex minus .2ex}{.8ex plus .2ex}
\pagestyle{fancy}
\fancyhf{}
\fancyhead[L]{\small\color{mutedink}When the secant numbers become periodic}
\fancyhead[R]{\small\color{mutedink}Research report}
\fancyfoot[C]{\small\thepage}
\renewcommand{\headrulewidth}{0.3pt}
\renewcommand{\footrulewidth}{0pt}
\setlength{\parskip}{0.32em}
\setlength{\parindent}{1.1em}
\setlength{\emergencystretch}{2em}
\setcounter{tocdepth}{2}
\numberwithin{equation}{section}
\newtheorem{theorem}{Theorem}[section]
\newtheorem{lemma}[theorem]{Lemma}
\newtheorem{proposition}[theorem]{Proposition}
\newtheorem{corollary}[theorem]{Corollary}
\theoremstyle{definition}
\newtheorem{definition}[theorem]{Definition}
\newtheorem{example}[theorem]{Example}
\newtheorem{remark}[theorem]{Remark}
\newcommand{\Z}{\mathbb Z}
\newcommand{\Q}{\mathbb Q}
\newcommand{\F}{\mathbb F}
\newcommand{\N}{\mathbb N}
\newcommand{\vp}{v_p}
\newcommand{\chiFour}{\chi_4}
\newcommand{\floor}[1]{\left\lfloor #1\right\rfloor}
\newcommand{\ceil}[1]{\left\lceil #1\right\rceil}
\DeclareMathOperator{\lcm}{lcm}
\DeclareMathOperator{\sech}{sech}
\DeclareMathOperator{\dens}{dens}
\newtcolorbox{resultbox}{colback=olivewash,colframe=ruleink,boxrule=.5pt,
  arc=1mm,left=3mm,right=3mm,top=2mm,bottom=2mm,breakable}
\lstset{basicstyle=\small\ttfamily,keywordstyle=\bfseries\color{oliveink},
  commentstyle=\color{mutedink},stringstyle=\color{oliveink},
  showstringspaces=false,breaklines=true,columns=fullflexible,
  keepspaces=true,frame=single,rulecolor=\color{ruleink},framerule=.4pt,
  xleftmargin=1mm,xrightmargin=1mm,aboveskip=1em,belowskip=1em}

\begin{document}
\thispagestyle{empty}
\begin{center}
{\small\sffamily\color{oliveink}INTEGER SEQUENCES \enspace / \enspace MODULAR ARITHMETIC}
\vspace{1.2em}

{\LARGE\bfseries When the Secant Numbers\\[.22em] Become Periodic}
\vspace{.8em}

{\large A counterexample, a complete classification,\\[.15em]
and a distribution law}
\vspace{1em}

{\small Prepared for Vladimir Reshetnikov \enspace\textbullet\enspace 19 September 2026}
\end{center}
\vspace{.4em}
\noindent\textcolor{ruleink}{\rule{\textwidth}{.6pt}}

\begin{abstract}
A conjecture recorded in OEIS A000364 asserts that the positive secant numbers,
starting at index one, are purely periodic modulo every positive integer $m$,
with a period dividing $\varphi(m)$. It is false: $a_1=1$ but
$a_{19}\equiv10\pmod{27}$, although $\varphi(27)=18$.
We replace the assertion with an exact theorem. Pure periodicity from index
one holds precisely when the odd part of $m$ is cube-free. For arbitrary $m$
we determine both the least eventual period and the exact first periodic
index. A finite alternating moment formula separates a periodic unit part
from a nilpotent part; the latter is exactly
$\chi_4(p)p^{2n}a_n$ modulo an odd prime power $p^e$.
A nonvanishing Fermat-quotient coefficient proves sharp period lifting,
and a falling-factorial argument handles powers of two.

The classification also yields the natural density of every bounded-preperiod
class, a finite-correction Euler product for those densities, and a limiting
mean first periodic index. In particular, the original assertion holds for
$8/(7\zeta(3))\approx95.0751\%$ of moduli. The mean is rigorously enclosed
near $1.05402581008837$. Exact Python computations, rational numerical
certificates, data, and a short standalone counterexample checker accompany
the report. The underlying Euler congruences are classical; the purpose is
to settle the specific recorded assertion, supply a complete specialized
account, and develop its distributional consequences, not to claim that
all constituent results are historically new.
\end{abstract}

\begin{resultbox}
\noindent\textbf{The corrected statement.} For every $m\ge1$, the secant
numbers are eventually periodic modulo $m$, with least eventual period
dividing $\varphi(m)$. Periodicity already holds for every $n\ge1$ if and
only if no odd prime cube divides $m$. The smallest exception is $m=27$.
\end{resultbox}

\clearpage
\begingroup
\small
\setlength{\parskip}{0pt}
\tableofcontents
\endgroup
\clearpage

\section{The problem, its status, and the first counterexample}
\label{sec:problem}

The positive secant numbers are defined by
\begin{equation}\label{eq:egf}
\sec x=\sum_{n\ge0}a_n\frac{x^{2n}}{(2n)!}.
\end{equation}
They begin
\[
1,\ 1,\ 5,\ 61,\ 1385,\ 50521,\ 2702765,\ 199360981,\ldots
\]
and form OEIS A000364. Combinatorially, $a_n$ counts alternating
permutations of $2n$ elements with a specified initial direction
\cite{oeis}. We use the signed Euler numbers only as an auxiliary notation:
\begin{equation}\label{eq:signed}
\sech t=\sum_{k\ge0}E_k\frac{t^k}{k!},\qquad
E_{2n}=(-1)^n a_n,\quad E_{2n+1}=0.
\end{equation}
This distinction matters: the full up/down sequence with generating
function $\sec x+\tan x$ is a different sequence, with a different indexing
and different periods.

\subsection{The assertion selected for investigation}

The OEIS entry, as consulted on 19 September 2026, still labels the following
assertion as a conjecture and attributes it to Peter Bala, 8 May 2023
\cite{oeis}:
\begin{equation}\label{eq:conjecture}
\text{For every }m\ge1,\quad (a_n\bmod m)_{n\ge1}
\text{ is purely periodic, with a period dividing }\varphi(m).
\end{equation}
Here $\varphi$ is Euler's totient function. The source's example is $m=21$,
whose least period is indeed $6$. Our status claim is deliberately precise:
this was a \emph{conjecture still displayed in the cited entry}. It does not
establish that no mathematician had previously noticed a counterexample or
that its resolution requires new congruence theory.

The distinction between \emph{pure} and \emph{eventual} periodicity is the
entire issue. The period-divisibility prediction survives; the claim that
there is no transient after index zero does not.

\subsection{A finite, independently checkable disproof}

From $\sec x\cos x=1$, coefficient comparison gives
\begin{equation}\label{eq:recurrence}
a_0=1,\qquad
a_n=\sum_{j=1}^{n}(-1)^{j+1}\binom{2n}{2j}a_{n-j}\quad(n\ge1).
\end{equation}
This integer recurrence gives
\[
a_{19}=23489580527043108252017828576198947741.
\]
Consequently
\begin{equation}\label{eq:counterexample}
a_1\equiv1\pmod{27},\qquad a_{19}\equiv10\pmod{27},
\qquad \varphi(27)=18.
\end{equation}
A period dividing $18$ would force equality at indices $1$ and $19$.
Thus \eqref{eq:counterexample} disproves \eqref{eq:conjecture}.
Appendix~\ref{app:checker} checks this using only the Python standard
library and integers of modest size.

Modulo $27$, the least eventual period is $18$ and the first periodic
index is $2$. The periodic block at indices $2$ through $19$ is
\[
\begin{split}
(&5,7,8,4,11,1,14,25,17,\\
 &22,20,19,23,16,26,13,2,10).
\end{split}
\]
At index $1$, the value is $1$, whereas periodic continuation backwards
from this block would demand $10$. This is a genuine transient, not a
mistaken guess about the length of a repeating block.

\subsection{What is classical, and what this report establishes}

Finite odd-power sums for Euler numbers and their Kummer-type congruences
are classical. The sum we use appears explicitly in Sun's paper
\cite[Eq.~(1.1)]{sun}; Mets\"ankyl\"a states the short Kummer congruence,
including the Euler factor responsible for the present counterexample
\cite[p.~161]{metsankyla}. Thus a short disproof of the OEIS assertion is
already available from established theory.

The modular arithmetic of the \emph{full} up/down sequence was studied by
Knuth and Buckholtz \cite{knuth}, then by Ramassamy \cite{ramassamy}.
A recent preprint of G\"ule\c{c} proves sharp odd-prime-power periods for
that full sequence and analyzes its preperiod \cite{gulec}. The present
report uses finite-moment and finite-frequency methods specialized to the
even subsequence. All proofs needed below are supplied explicitly.
We do not claim to settle the additional preperiod conjectures in that
preprint.

The deliverable is a complete resolution of \eqref{eq:conjecture} as
written, together with exact local and global invariants, distribution
formulas, and reproducible computations. The density results are deductions
proved here; an exhaustive priority claim for them is not made.

\section{The complete classification}
\label{sec:main}

\begin{definition}
For a fixed modulus $m\ge1$, let $d(m)$ be the least positive eventual
period of $(a_n\bmod m)_{n\ge0}$. Let $s(m)$ be the least $s\ge0$ such that
\[
a_{n+d(m)}\equiv a_n\pmod m\qquad(n\ge s).
\]
We call $s(m)$ the \emph{first periodic index}, or \emph{preperiod length}
when the sequence is indexed from zero. Thus the assertion concerning
indices $n\ge1$ is equivalent to $s(m)\le1$.
\end{definition}

For an odd prime $p$, put
\begin{equation}\label{eq:P}
P_e(p)=p^{e-1}\lcm\left(2,\frac{p-1}{2}\right)
=\begin{cases}
 p^{e-1}(p-1)/2,&p\equiv1\pmod4,\\
 p^{e-1}(p-1),&p\equiv3\pmod4.
\end{cases}
\end{equation}
We write $v_p(u)$ for the exponent of $p$ in a nonzero integer $u$, and
$\chi_4(p)=(-1)^{(p-1)/2}$.

\begin{theorem}[Exact local and global invariants]\label{thm:main}
Every modulus has an eventually periodic secant sequence. For an odd
prime $p$ and $e\ge1$,
\begin{align}
 d(p^e)&=P_e(p),\label{eq:oddperiod}\\
 s(p^e)&=1+\max\left\{0\le n<\ceil{e/2}:
                  2n+v_p(a_n)<e\right\}.\label{eq:oddpreperiod}
\end{align}
The maximum is over a nonempty set. For powers of two,
\begin{equation}\label{eq:dyadic}
s(2^b)=0,\qquad
d(2^b)=\begin{cases}1,&b=1,2,\\2^{b-1},&b\ge3.\end{cases}
\end{equation}
For $m=1$, set $s(1)=0$ and $d(1)=1$. For every other modulus,
\begin{equation}\label{eq:CRT}
s(m)=\max_{p^e\parallel m}s(p^e),\qquad
d(m)=\lcm_{p^e\parallel m}d(p^e).
\end{equation}
In particular, $d(m)\mid\varphi(m)$ for all $m\ge1$.
\end{theorem}

The proof is divided into a finite-transient lemma, the odd-prime argument,
the dyadic argument, and the Chinese remainder theorem. Formula
\eqref{eq:oddpreperiod} is finite and effective: it needs only the first
$\ceil{e/2}$ secant numbers, no matter how large $p^e$ or its period is.

\begin{corollary}[Exact replacement of the conjecture]\label{cor:cubefree}
For every $m\ge1$, the following are equivalent:
\begin{enumerate}[label=(\roman*),nosep]
\item $(a_n\bmod m)_{n\ge1}$ is purely periodic;
\item $s(m)\le1$;
\item no odd prime cube divides $m$.
\end{enumerate}
When these conditions hold, the least period divides $\varphi(m)$.
The smallest modulus violating the original assertion is $27$.
\end{corollary}

\begin{proof}[Deduction from the theorem]
If every odd prime exponent is at most two, then for $n\ge1$ one has
$2n\ge e$, so the set in \eqref{eq:oddpreperiod} consists only of $n=0$.
Thus every odd local preperiod is $1$; powers of two cause no transient.
Conversely, if $e\ge3$, then $a_1=1$ gives
$2+v_p(a_1)=2<e$, so $s(p^e)\ge2$.
The least positive odd prime cube is $3^3=27$.
\end{proof}

\begin{table}[htbp]
\centering
\caption{Exact invariants; indexing starts at zero.}
\label{tab:examples}
\begin{tabular}{r r r @{
\qquad} r r r}
\toprule
$m$ & $s(m)$ & $d(m)$ & $m$ & $s(m)$ & $d(m)$\\
\midrule
1&0&1&27&2&18\\
4&0&1&32&0&16\\
8&0&4&54&2&18\\
9&1&6&81&2&54\\
16&0&8&125&2&50\\
21&1&6&243&3&162\\
25&1&10&625&2&250\\
&&&3125&2&1250\\
\bottomrule
\end{tabular}
\end{table}

Notice the two different sources of arithmetic behavior. The period is
an explicit multiplicative-order quantity, assembled by least common
multiples. The first periodic index depends on divisibility of individual
Euler numbers. For example, $s(5^5)=2$, not the generic upper bound $3$,
because $a_2=5$ supplies an extra factor of $5$.

\section{A periodic sequence plus a finite transient}
\label{sec:transient}

The following elementary observation prevents an important logical
mistake: observing that one proposed period fails at an early index is not,
by itself, proof that \emph{every} possible period fails there.

\begin{lemma}[Finite-transient separation]\label{lem:transient}
Let $A$ be an abelian group. Suppose a sequence in $A$ has the form
$x_n=u_n+b_n$, where $u$ is purely periodic from index zero and $b_n=0$
for all sufficiently large $n$. Then:
\begin{enumerate}[label=(\roman*),nosep]
\item the eventual periods of $x$ are exactly the periods of $u$;
\item the least eventual period of $x$ is the least period of $u$;
\item if $b$ is nonzero somewhere, the first periodic index of $x$ is
$1+\max\{n:b_n\ne0\}$; if $b$ is identically zero, it is $0$.
\end{enumerate}
Every eventual period starts working at that same first periodic index.
\end{lemma}

\begin{proof}
Let $L$ be any pure period of $u$. If $T$ is an eventual period of $x$,
then it is an eventual period of $u$, because $x=u$ on a sufficiently late
tail. For any fixed $n\ge0$, choose $k$ so large that $n+kL$ lies in that
tail. Pure $L$-periodicity gives
\[
u_{n+T}=u_{n+T+kL}=u_{n+kL}=u_n.
\]
Thus $T$ is a pure period of $u$. Conversely, every period of $u$ is
clearly an eventual period of $x$.

Suppose now that $T$ works for $x$ from index $s$. Since it works for $u$
everywhere, subtraction shows $b_{n+T}=b_n$ for all $n\ge s$.
Iterating this equality into the zero tail gives $b_n=0$ for every $n\ge s$.
The converse is immediate. This proves the exact onset statement.
\end{proof}

We will also use the usual divisibility property of periods. If a purely
periodic one-sided sequence has periods $P>Q$, then $P-Q$ is a period:
\[
u_{n+P-Q}=u_{n+P}=u_n.
\]
The middle equality uses the period $Q$ at the nonnegative index
$n+P-Q$. The Euclidean algorithm shows that the least period divides every
period. The same conclusion holds for eventual periods by passing to a
tail. In particular, reducing a sequence modulo a divisor of its modulus
forces its least eventual period to divide every eventual period at the
larger modulus.

\section{Finite moments and the exact odd-prime transient}
\label{sec:moments}

\subsection{An integer identity before reduction}

\begin{lemma}[Finite alternating moment formula]\label{lem:moment}
For every positive odd integer $q$ and every $k\ge0$,
\begin{equation}\label{eq:momentE}
E_k\equiv\sum_{j=0}^{q-1}(-1)^j(2j+1)^k\pmod q.
\end{equation}
Consequently, for $n\ge0$,
\begin{equation}\label{eq:momentA}
a_n\equiv\sum_{j=0}^{q-1}(-1)^j\bigl(-(2j+1)^2\bigr)^n\pmod q.
\end{equation}
At $n=0$, the convention $x^0=1$ includes residue classes $x=0$.
\end{lemma}

\begin{proof}
Work first in $\Q[[t]]$, not in a ring where factorials might be
noninvertible. A finite geometric sum yields
\[
\sum_{j=0}^{q-1}(-1)^j e^{(2j+1)t}
=\frac{e^t(1+e^{2qt})}{1+e^{2t}}
=\frac{1+e^{2qt}}{2\cosh t}.
\]
Comparing the coefficients of $t^k/k!$ gives the exact identity
\[
\sum_{j=0}^{q-1}(-1)^j(2j+1)^k
=E_k+\frac12\sum_{\ell=1}^k\binom{k}{\ell}(2q)^\ell E_{k-\ell}.
\]
Every term in the sum on the right is an integer multiple of $q$.
The signed Euler numbers are integers by the coefficient recurrence for
$\sech t\cosh t=1$. Reduction proves \eqref{eq:momentE}; substituting
$k=2n$ and multiplying by $(-1)^n$ proves \eqref{eq:momentA}.
\end{proof}

The congruence is known \cite[Eq.~(1.1)]{sun}; the proof above records
explicitly why no division by a factorial modulo $q$ is involved.

\subsection{Unit and nonunit contributions}

Fix an odd prime power $q=p^e$. In $\Z/p^e\Z$, split the last sum into
\begin{align*}
 U_n&=\sum_{\substack{0\le j<p^e\\p\nmid2j+1}}
        (-1)^j\bigl(-(2j+1)^2\bigr)^n,\\
 B_n&=\sum_{\substack{0\le j<p^e\\p\mid2j+1}}
        (-1)^j\bigl(-(2j+1)^2\bigr)^n.
\end{align*}
Then $a_n=U_n+B_n$ in this residue ring. Every ratio in $U_n$ is a unit;
every ratio in $B_n$ is divisible by $p^2$.

\begin{proposition}[Exact nonunit contribution]\label{prop:depletion}
For every $n\ge0$,
\begin{equation}\label{eq:depletion}
 B_n=\chi_4(p)p^{2n}a_n,\qquad
 U_n=(1-\chi_4(p)p^{2n})a_n\quad\text{in }\Z/p^e\Z.
\end{equation}
Moreover $U$ is purely $P_e(p)$-periodic and $B_n=0$ whenever $2n\ge e$.
\end{proposition}

\begin{proof}
The odd multiples of $p$ in the required range are exactly
\[
2j+1=p(2r+1),\quad
j=pr+(p-1)/2,\quad0\le r<p^{e-1}.
\]
Since $p$ is odd, $(-1)^j=\chi_4(p)(-1)^r$. Therefore
\[
 B_n=\chi_4(p)p^{2n}
 \sum_{r=0}^{p^{e-1}-1}(-1)^r\bigl(-(2r+1)^2\bigr)^n.
\]
For $n\ge1$, Lemma~\ref{lem:moment} makes the inner sum congruent to $a_n$
modulo $p^{e-1}$. Multiplication by $p^{2n}$ raises the precision to at
least $p^e$. This also covers $e=1$, where the stated inner congruence is
modulo $1$. For $n=0$, both the inner alternating sum and $a_0$ are
exactly $1$, so the identity holds directly.

The integer $P_e(p)$ is even and satisfies
$\varphi(p^e)\mid2P_e(p)$. For a unit $x$ modulo $p^e$, Euler's theorem
therefore gives
$(-x^2)^{P_e(p)}=1$. This proves pure periodicity of $U$.
Finally every term in $B_n$ is divisible by $p^{2n}$.
\end{proof}

Apply Lemma~\ref{lem:transient}. The nonzero indices of $B$ are precisely
\[
2n+v_p(a_n)<e.
\]
They lie below $\ceil{e/2}$ and include $n=0$. This proves
\eqref{eq:oddpreperiod} and the upper bound $d(p^e)\mid P_e(p)$.
Only sharpness of the period remains.

The factor $1-\chi_4(p)p^{2n}$ is not an arbitrary correction introduced to
fit the counterexample. It is exactly what remains after discarding the
nonunit bases in the finite moment sum. Its appearance is also the
classical Euler factor in Kummer congruences \cite{metsankyla}.

\section{Why the odd-prime-power period is minimal}
\label{sec:oddperiod}

\subsection{Distinct exponential sequences are independent}

\begin{lemma}[Finite-field independence]\label{lem:vandermonde}
Let $\lambda_1,\ldots,\lambda_r$ be distinct nonzero elements of a field.
If
\[
\sum_{i=1}^r c_i\lambda_i^n=0
\]
holds for $r$ consecutive integer indices $n=N,\ldots,N+r-1$, where
$N\ge0$, then every $c_i=0$.
\end{lemma}

\begin{proof}
The coefficient matrix is a Vandermonde matrix multiplied by the diagonal
matrix with entries $\lambda_i^N$. Its determinant is
\[
\left(\prod_i\lambda_i^N\right)
\prod_{i<j}(\lambda_j-\lambda_i),
\]
which is nonzero.
\end{proof}

\subsection{The period modulo an odd prime}

In the moment sum modulo $p$, pair $j$ with $p-1-j$. The corresponding
odd residues are negatives of one another, and the two signs $(-1)^j$
agree. Apart from the unique zero residue, the negative squares
$\lambda=-(2j+1)^2$ therefore occur in pairs, each with coefficient $+2$
or $-2$. These coefficients are nonzero in $\F_p$.

For positive indices the zero-residue term is absent. Hence an eventual
period $T$ gives
\[
\sum_{\lambda}c_\lambda(\lambda^T-1)\lambda^n=0
\]
on a tail. By Lemma~\ref{lem:vandermonde}, every $\lambda^T=1$.
Taking $\lambda=-1$ forces $T$ to be even. Then $(-x^2)^T=1$ for all
$x\in\F_p^\times$ implies $x^{2T}=1$ for all such $x$; since the
multiplicative group of a finite field is cyclic,
\[
p-1\mid2T.
\]
The least positive $T$ satisfying these two conditions is
$\lcm(2,(p-1)/2)=P_1(p)$. The upper bound already proved shows
\begin{equation}\label{eq:baseperiod}
d(p)=P_1(p).
\end{equation}
This is the classical odd-prime period formula recorded in the OEIS with
reference to Knuth and Buckholtz \cite{oeis,knuth}.

\subsection{The lifting step and its nonzero coefficient}

Suppose by induction that $d(p^{e-1})=P_{e-1}(p)$, where $e\ge2$.
Reduction modulo $p^{e-1}$ and the upper bound imply
\[
P_{e-1}(p)\mid d(p^e)\mid P_e(p)=pP_{e-1}(p).
\]
There are only two possibilities. Set $T=P_{e-1}(p)$; it suffices to
prove that $T$ is not an eventual period modulo $p^e$.

Put
\[
\delta=\frac{2P_1(p)}{p-1}\in\{1,2\},\qquad
q_p(x)=\frac{x^{p-1}-1}{p}\quad(p\nmid x).
\]
The integer $q_p(x)$ is the Fermat quotient. Elementary binomial expansion
gives
\begin{equation}\label{eq:lift}
x^{2T}=(1+pq_p(x))^{\delta p^{e-2}}
\equiv1+\delta p^{e-1}q_p(x)\pmod{p^e}.
\end{equation}
For completeness, the familiar step
$(1+py)^{p^k}\equiv1+p^{k+1}y\pmod{p^{k+2}}$
follows by induction on $k$, expanding one additional $p$th power.
Raising to $\delta=1$ or $2$ gives \eqref{eq:lift}.

For $n$ large enough that the nonunit terms vanish, the moment formula and
$T$ even yield
\begin{equation}\label{eq:normalizeddefect}
\frac{a_{n+T}-a_n}{p^{e-1}}
\equiv\delta
\sum_{\substack{0\le j<p^e\\p\nmid2j+1}}
(-1)^j q_p(2j+1)\bigl(-(2j+1)^2\bigr)^n\pmod p.
\end{equation}
The quotient on the left is an integer on a sufficiently late tail by
the inductive period statement. Equivalently, one may first form the
congruence modulo $p^e$ with a factor $p^{e-1}$ on both sides, then reduce
the quotient modulo $p$.

Group the right side by the distinct negative squares modulo $p$.
We compute just one coefficient: the coefficient of $\lambda=-1$,
coming from $x\equiv\pm1\pmod p$. Let $M=p^{e-1}$. The two groups are
\[
\begin{array}{c|c|c}
 x & j & q_p(x)\pmod p\\ \hline
 1+2p\ell & p\ell & -2\ell\\
 -1+2p(\ell+1) & p-1+p\ell & 2(\ell+1)
\end{array}
\qquad(0\le\ell<M).
\]
The quotient formulas follow by expanding $x^{p-1}$ modulo $p^2$.
Both rows have sign $(-1)^j=(-1)^\ell$. Since $M$ is odd and $p\mid M$,
\[
\sum_{\ell=0}^{M-1}(-1)^\ell=1,\qquad
\sum_{\ell=0}^{M-1}(-1)^\ell\ell=\frac{M-1}{2}.
\]
The first row therefore contributes $-(M-1)\equiv1\pmod p$, and the
second contributes $M+1\equiv1\pmod p$. The coefficient of the
$(-1)^n$ mode in \eqref{eq:normalizeddefect} is consequently
\begin{equation}\label{eq:nonzero}
2\delta\not\equiv0\pmod p.
\end{equation}
By Lemma~\ref{lem:vandermonde}, the normalized defect cannot vanish for all
sufficiently large $n$. Hence $T$ is not an eventual period modulo $p^e$.
The other possibility is forced:
\[
d(p^e)=pP_{e-1}(p)=P_e(p).
\]
This completes the proof of \eqref{eq:oddperiod}.

The important feature is that the obstruction to losing a factor $p$
is the fixed coefficient $2\delta$. It does not depend on the
nonvanishing of a particular large Euler number, and it cannot disappear
at an exceptional prime. The use of grouped finite frequencies is closely
related to the method used for the full up/down sequence in \cite{gulec};
here the even subsequence allows the argument to stay in $\F_p$.

\section{Powers of two: a self-contained Stern argument}
\label{sec:two}

Odd finite moments do not directly address dyadic moduli. We give a short
proof of the exact signed-Euler congruence traditionally associated with
Stern, then translate it carefully to positive secant numbers. For its
classical provenance and other proofs, see Sun \cite{sun}.

\begin{lemma}[Dyadic isometry on even Euler indices]\label{lem:stern}
If $N,M$ are distinct nonnegative even integers, then
\begin{equation}\label{eq:stern}
v_2(E_N-E_M)=v_2(N-M).
\end{equation}
Moreover, for even $N\ge0$,
\begin{equation}\label{eq:Emod8}
E_N\equiv1-N\pmod8.
\end{equation}
\end{lemma}

\begin{proof}
Let $R=\Z_{(2)}$, the ring of rational numbers with odd denominator.
The formal series
\[
u(t)=\frac{e^{2t}-1}{2}
=\sum_{j\ge1}\frac{2^{j-1}}{j!}t^j
\]
has coefficients in $R$, since $v_2(j!)\le j-1$. Its constant term is
zero, so
\[
\tanh t=\frac{u(t)}{1+u(t)}\in R[[t]].
\]
Write $\tanh t=\sum_{j\ge1,\ j\text{ odd}}b_jt^j$, where $b_j\in R$
and $b_1=1$.
The exact identity $\sech t=e^t(1-\tanh t)$ gives
\begin{equation}\label{eq:falling}
E_N=1-\sum_{\substack{1\le j\le N\\j\text{ odd}}}
b_j(N)_{\underline j},
\qquad (X)_{\underline j}=X(X-1)\cdots(X-j+1).
\end{equation}
For comparison at $N$ and $M$, use the same finite summation range up to
$\max(N,M)$: terms beyond either nonnegative integer index vanish.

For odd $j\ge3$, the product $(X)_{\underline j}$ has at least two
factors with even offsets, namely $X$ and $X-2$. Telescope the difference
of the products evaluated at $N$ and $M$, changing one factor at a time.
Each resulting summand contains one factor $N-M$. After the changed
factor is omitted, at least one even-offset factor remains; its value at
either even argument is even. Hence
\[
(N)_{\underline j}-(M)_{\underline j}\in2(N-M)\Z
\qquad(j\ge3\text{ odd}).
\]
The $j=1$ term in \eqref{eq:falling} is $-N$. It follows that
\[
E_N-E_M=-(N-M)+2(N-M)r\qquad\text{for some }r\in R.
\]
The factor $-1+2r$ is a unit in $R$, proving \eqref{eq:stern}.

Finally, for even $N$ and odd $j\ge3$, the two factors $N$ and $N-2$
include a multiple of $4$ and another multiple of $2$. Thus
$(N)_{\underline j}\in8\Z$. The same conclusion is automatic for a
zero product. In \eqref{eq:falling}, all terms except $1-N$ vanish
modulo $8$, proving \eqref{eq:Emod8}.
\end{proof}

By $a_n=(-1)^nE_{2n}$,
\begin{equation}\label{eq:amod8}
a_n\equiv(-1)^n(1-2n)\pmod8.
\end{equation}
The repeating residue block is $(1,1,5,5)$, starting at $n=0$.
In particular, every $a_n\equiv1\pmod4$, so the period modulo $2$ and
$4$ is $1$.

For $b\ge3$, every eventual period must be divisible by $4$, since the
block modulo $8$ has least period $4$. For an even shift $T$, signs at
$n$ and $n+T$ agree, and \eqref{eq:stern} gives, at every $n\ge0$,
\[
v_2(a_{n+T}-a_n)=v_2(2T)=1+v_2(T).
\]
Thus periodicity modulo $2^b$ is equivalent, among the possible shifts,
to $2^{b-1}\mid T$. The least such shift is $2^{b-1}$, and it works
already at zero. This proves \eqref{eq:dyadic}.

\section{Global assembly and a formula for every failure}
\label{sec:global}

\subsection{Chinese remainders and the exact onset}

A congruence modulo $m$ is equivalent to the congruence at every exact
prime-power factor of $m$. An eventual period must therefore be a
multiple of each local least period, and the least common multiple is
itself an eventual period. This proves the period formula in
\eqref{eq:CRT}.

For the onset, each odd local sequence is a periodic sequence plus a
finite transient, as in Lemma~\ref{lem:transient}; the dyadic local sequence
has no transient. Every local eventual period therefore has the same
least valid onset. A global period cannot begin earlier than any local
one, while the least common multiple works after the latest local onset.
This proves the other formula in \eqref{eq:CRT}.

Every local period divides $\varphi(p^e)$. Since
$\varphi(m)=\prod_{p^e\parallel m}\varphi(p^e)$, every local period
divides $\varphi(m)$, and so does their least common multiple.
This finishes the proof of Theorem~\ref{thm:main}.

\subsection{The defect at a totient shift}

\begin{proposition}[All-index defect]\label{prop:defect}
Let $p^e$ be an odd prime-power divisor of $m$. For every $n\ge0$,
\begin{equation}\label{eq:defect}
a_{n+\varphi(m)}-a_n
\equiv-\chi_4(p)p^{2n}a_n\pmod{p^e}.
\end{equation}
In particular,
\begin{equation}\label{eq:firstdefect}
a_{1+\varphi(m)}-1\equiv-\chi_4(p)p^2\pmod{p^e}.
\end{equation}
\end{proposition}

\begin{proof}
The shift $\varphi(m)$ is a multiple of $P_e(p)$, so
$U_{n+\varphi(m)}=U_n$. Also $\varphi(m)\ge\varphi(p^e)\ge e$ for
odd $p$, so $B_{n+\varphi(m)}=0$. Since $a=U+B$, the difference is
$-B_n$. Substitute \eqref{eq:depletion} and then $a_1=1$.
\end{proof}

Consequently every modulus excluded by the cube-free criterion fails
at the very first tested index for the conjectured totient shift. It is
not necessary to search far into a period for a mismatch.
For the infinite family $m=p^3$,
\begin{equation}\label{eq:infinitefamily}
a_{1+p^2(p-1)}\equiv1-\chi_4(p)p^2\pmod{p^3}.
\end{equation}
For $p=3,5,7$, the right sides are respectively $10,101,50$ modulo
$27,125,343$.

This also shows why replacing the actual least period by a much larger
multiple cannot remove the transient: the unit contribution stays fixed,
while the nonunit contribution at the early index is still present.

\subsection{Two useful local examples}

Modulo $3$, \eqref{eq:momentA} has two equal unit modes. For $n\ge1$,
\begin{equation}\label{eq:mod3}
a_n\equiv2(-1)^n\pmod3.
\end{equation}
Thus $3\nmid a_n$ for every $n\ge1$, and
\[
s(3^e)=\ceil{e/2}\qquad(e\ge1).
\]
The generic upper bound is attained at every power of $3$.

At $p=5,e=5$, the candidate transient indices are $0,1,2$.
Their values of $2n+v_5(a_n)$ are $0,2,5$, respectively.
Only $0$ and $1$ are nonzero-transient indices modulo $5^5$, so
$s(3125)=2$, whereas $d(3125)=1250$.
Thus an exact preperiod theorem must retain divisibility information;
$s(p^e)=\ceil{e/2}$ is not universally valid.

\section{Generating functions and effective computation}
\label{sec:algorithms}

\subsection{A rational ordinary generating function modulo an odd integer}

The ordinary generating series is used formally, without asserting
analytic convergence over the complex numbers. Define
\[
G_q(z)=\sum_{n\ge0}(a_n\bmod q)z^n\in(\Z/q\Z)[[z]].
\]
For odd $q$, summing \eqref{eq:momentA} coefficientwise gives
\begin{equation}\label{eq:modgf}
G_q(z)=\sum_{j=0}^{q-1}\frac{(-1)^j}{1+(2j+1)^2z}.
\end{equation}
Every denominator has constant coefficient $1$, hence is invertible as
a formal power series even when $q$ is composite.

At a prime power, denominators with a nonunit ratio contribute only
polynomials: if $p\mid2j+1$, then the geometric series stops after at
most $\ceil{e/2}$ terms because its ratio is nilpotent. Unit ratios
produce the periodic part. Thus the polynomial-plus-periodic split of
the coefficient sequence is visible directly in the rational series.

For an arbitrary modulus, put $s=s(m)$, $d=d(m)$ and write
$\bar a_n=a_n\bmod m$. Then
\begin{equation}\label{eq:generalGF}
G_m(z)=\sum_{n=0}^{s-1}\bar a_nz^n
 +\frac{z^s}{1-z^d}\sum_{j=0}^{d-1}\bar a_{s+j}z^j.
\end{equation}
An empty initial sum is zero. We claim this explicit representation, not
an irreducible denominator or a field-like notion of minimal denominator
degree over a residue ring with zero divisors.

\subsection{Computing periods without enumerating a period}

Given a factorization of $m$, the period is obtained by integer
exponentiation and least common multiples. To compute the onset, let
$H$ be the largest $\ceil{e/2}$ among odd prime powers dividing $m$.
Generate $a_0,\ldots,a_{H-1}$ and inspect the valuations in
\eqref{eq:oddpreperiod}. The number of coefficient operations in the
binomial recurrence is $O(H^2)$ when binomial coefficients are available;
this is an arithmetic-operation estimate, not a unit-cost bit complexity
claim for arbitrarily large integers.

The supplied implementation uses trial division for factorization and
primality checks. It is intentionally transparent, not a fast
factorization package. It is suitable for modest moduli and for
prime-power experiments with modest primes. Factoring a huge arbitrary
modulus can be the dominant cost and is not hidden in the period formula.

For a huge index $n$ and a modest odd modulus $q$,
\eqref{eq:momentA} computes $a_n\bmod q$ using $q$ modular
exponentiations, or $O(q\log(n+1))$ modular multiplications with binary
exponentiation. Extra working storage apart from the input integers is
constant. If $s(q)$ and $d(q)$ are known and $n\ge s(q)$, the index may
instead be reduced to
\[
s(q)+(n-s(q))\bmod d(q).
\]
These are complementary approaches: the moment formula handles huge
indices without ever constructing huge secant numbers, while the onset
formula uses only a short initial segment.

\subsection{Independent coefficient generation}

For testing, the archive also constructs the coefficients through the
Entringer triangle, rather than reusing \eqref{eq:recurrence}. Set
$T(0,0)=1$, and for $r\ge1$ use
\[
T(r,0)=0,\qquad
T(r,k)=T(r,k-1)+T(r-1,r-k)\quad(1\le k\le r).
\]
Then $T(2n,2n)=a_n$. This is a standard alternating-permutation triangle;
its row-end construction is also described in the literature on the
full up/down numbers \cite{ramassamy}. The implementation stores only
the current row and the required even-row endpoints. Comparison of two
different constructions helps detect indexing, sign, and recurrence
implementation errors; it is not a substitute for the proofs above.

\section{How often is the original assertion true?}
\label{sec:density}

The counterexamples have an exact arithmetic description, so we can count
them rather than just exhibit them. More generally, we can determine the
natural density of moduli with any prescribed upper bound on their first
periodic index.

For a set $S$ of positive integers, its natural density, when it exists, is
\[
\dens(S)=\lim_{X\to\infty}\frac{\#\{m\le X:m\in S\}}{\floor X}.
\]
For $N\ge1$ and an odd prime $p$, define
\begin{equation}\label{eq:cutoff}
c_p(N)=\min_{n\ge N}\bigl(2n+v_p(a_n)\bigr).
\end{equation}
This looks like an infinite minimization, but is actually a finite one.

\begin{lemma}[Finite cutoff computation]\label{lem:cutoff}
Let $v=v_p(a_N)$. Then
\begin{equation}\label{eq:finitecutoff}
c_p(N)=\min_{N\le n\le N+\floor{v/2}}
             \bigl(2n+v_p(a_n)\bigr).
\end{equation}
In particular, $c_p(N)\ge2N$, with equality whenever $p\nmid a_N$.
\end{lemma}

\begin{proof}
The term at $n=N$ gives the upper bound $2N+v$. If
$n>N+\floor{v/2}$, then $2n>2N+v$, so that index cannot minimize the
expression. The final assertions follow at once.
\end{proof}

\begin{theorem}[Distribution of the first periodic index]\label{thm:density}
For $N\ge1$, the set $\{m:s(m)\le N\}$ has natural density
\begin{equation}\label{eq:densityproduct}
D_N=\prod_{p\text{ odd}}\left(1-p^{-(c_p(N)+1)}\right).
\end{equation}
Equivalently,
\begin{equation}\label{eq:densityzeta}
D_N=\frac{1}{(1-2^{-(2N+1)})\zeta(2N+1)}
\prod_{\substack{p\mid a_N\\p\text{ odd}}}
\frac{1-p^{-(c_p(N)+1)}}{1-p^{-(2N+1)}}.
\end{equation}
The correction product is finite. More precisely, for fixed $N$,
\begin{equation}\label{eq:countasymptotic}
\#\{m\le X:s(m)\le N\}
=D_N X+O_N\!\left(X^{1/(2N+1)}\right).
\end{equation}
\end{theorem}

\begin{proof}
By the exact transient criterion, $s(m)\le N$ if and only if, at every
odd prime $p$, all nonunit contributions vanish at every index $n\ge N$.
Writing $e=v_p(m)$, this is equivalent to
\[
e\le2n+v_p(a_n)\quad\text{for every }n\ge N,
\]
or $v_p(m)\le c_p(N)$. There is no restriction at $2$.
Thus the set is exactly the set of integers avoiding divisibility by
each $p^{c_p(N)+1}$ for odd primes $p$.

Here is a counting proof that also justifies the infinite Euler product.
Set $r=2N+1$ and let $S$ be the finite set of odd prime divisors of $a_N$.
Outside $S$, the forbidden exponent is exactly $r$. For $p\in S$, put
$b_p=c_p(N)+1\ge r$. Inclusion--exclusion at the exceptional primes,
and M\"obius inversion elsewhere, give the exact count
\[
\sum_{J\subseteq S}(-1)^{|J|}
\sum_{\substack{d\ge1\\(d,\,2\prod_{p\in S}p)=1}}
\mu(d)\left\lfloor\frac{X}{d^r C_J}\right\rfloor,
\qquad C_J=\prod_{p\in J}p^{b_p}.
\]
Only $d\le(X/C_J)^{1/r}$ contribute. Replacing the floor by its
argument costs $O_N(X^{1/r})$, since there are that many potentially
nonzero terms for each of finitely many $J$. Extending the main-term
sum over $d$ to infinity also costs $O_N(X^{1/r})$, using
$\sum_{d>Y}d^{-r}=O_r(Y^{1-r})$.

The coefficient of $X$ is
\[
\prod_{p\in S}(1-p^{-b_p})
\prod_{\substack{p\text{ odd}\\p\notin S}}(1-p^{-r}).
\]
This proves \eqref{eq:densityproduct} and \eqref{eq:countasymptotic}.
The Euler product for $\zeta(r)$, valid since $r>1$, gives
\eqref{eq:densityzeta}. Absolute convergence also follows from
$c_p(N)+1\ge r\ge3$.
\end{proof}

\begin{corollary}[Density of the moduli satisfying the conjecture]
The natural density of moduli for which \eqref{eq:conjecture} is true is
\begin{equation}\label{eq:cubefreedensity}
D_1=\prod_{p\text{ odd}}(1-p^{-3})
=\frac{8}{7\zeta(3)}
=0.95075128294937996420\ldots.
\end{equation}
The counterexample moduli have density $1-D_1$, approximately $4.92487\%$.
\end{corollary}

\begin{proof}
Since $a_1=1$, there are no exceptional primes in
\eqref{eq:densityzeta}; all cutoffs are $c_p(1)=2$.
\end{proof}

The restriction is that the \emph{odd part} be cube-free; arbitrary powers
of $2$ are allowed. Accordingly the density is not $1/\zeta(3)$.
The factor $(1-2^{-3})^{-1}=8/7$ accounts exactly for the unrestricted
power of two.

\section{The whole distribution and its asymptotic tail}
\label{sec:tail}

\subsection{The first exceptional correction}

At $N=2$, $a_2=5$, so $c_5(2)=5$ and all other odd-prime cutoffs are $4$.
Therefore
\begin{equation}\label{eq:D2}
D_2=\frac{1-5^{-6}}{1-5^{-5}}
\frac{1}{(1-2^{-5})\zeta(5)}.
\end{equation}
The corresponding moduli have odd prime exponents at most $4$, except
that exponent $5$ is allowed at the prime $5$. This is the density-level
reflection of $s(5^5)=2$.

At $N=10$, one must not replace every exceptional cutoff by $2N+1$:
\[
a_{10}=5^2\cdot41737\cdot354957173.
\]
Here $c_5(10)=22$, while the other two cutoffs are $21$. The finite
minimization in Lemma~\ref{lem:cutoff} captures such higher-valuation effects
automatically.

\begin{table}[htbp]
\centering
\caption{Densities $D_N=\dens\{m:s(m)\le N\}$, rounded to 18 decimal
places. Full rational enclosures are supplied in the archive.}
\label{tab:densities}
\begin{tabular}{r l r l}
\toprule
$N$ & $D_N$ & $N$ & $D_N$\\
\midrule
1&0.950751282949379964&6&0.999999372600335288\\
2&0.995751538184580342&7&0.999999930275301732\\
3&0.999528673601264576&8&0.999999992256209180\\
4&0.999949067031528573&9&0.999999999139555886\\
5&0.999994333981021976&10&0.999999999904400848\\
\bottomrule
\end{tabular}
\end{table}

Only powers of two have $s(m)=0$, and their density is zero. Define $D_0=0$.
For every $N\ge1$, the moduli with exact first periodic index $N$ have
density $D_N-D_{N-1}$. Thus the cumulative densities determine an entire
probability distribution in the limit of uniformly sampled moduli.

\subsection{The tail is controlled by the prime three}

\begin{theorem}[Asymptotic distribution tail]\label{thm:tail}
As $N\to\infty$,
\begin{equation}\label{eq:tail}
1-D_N=3^{-(2N+1)}+O\!\left(5^{-(2N+1)}\right).
\end{equation}
In particular, the limiting distribution has finite moments of every
positive integer order.
\end{theorem}

\begin{proof}
Equation \eqref{eq:mod3} shows that $c_3(N)=2N$ for every $N\ge1$.
Thus the set of multiples of $3^{2N+1}$ is included in the excluded
moduli and has density $3^{-(2N+1)}$. Every other prime has forbidden
exponent at least $2N+1$. A union bound gives, with $r=2N+1$,
\[
3^{-r}\le1-D_N\le3^{-r}+\sum_{p\ge5}p^{-r}.
\]
For $r\ge3$, the last sum is at most
\[
\sum_{k=5}^\infty k^{-r}
\le5^{-r}+\frac{5^{1-r}}{r-1}=O(5^{-r}),
\]
with an absolute constant. This proves the expansion. Exponential decay
of the tail implies finiteness of every polynomial moment by the usual
discrete tail-sum formula.
\end{proof}

The theorem has a concrete interpretation: among very long transients,
the overwhelmingly dominant cause is a high power of $3$ in the modulus.
Exceptional divisibility of secant numbers can only shorten transients;
it cannot create a forbidden exponent below $2N+1$ at another prime.

\section{The mean first periodic index}
\label{sec:mean}

\begin{theorem}[Mean and a certified enclosure]\label{thm:mean}
The arithmetic mean of $s(m)$ over initial intervals has a limit:
\begin{equation}\label{eq:mean}
\mathcal M=\lim_{X\to\infty}\frac1{\floor X}\sum_{m\le X}s(m)
=1+\sum_{N\ge1}(1-D_N).
\end{equation}
It satisfies the rigorous bounds
\begin{equation}\label{eq:meaninterval}
1.05402581008837151<\mathcal M<1.05402581008837162.
\end{equation}
\end{theorem}

\begin{proof}[Proof of the limit formula]
For every nonnegative integer-valued function,
\[
s(m)=\sum_{N\ge0}\mathbf1_{\{s(m)>N\}}.
\]
For $N=0$, the limiting proportion is $1$, since $s(m)=0$ only at powers
of two. For fixed $N\ge1$, the limiting proportion is $1-D_N$.
We must justify passing the limit through the infinite sum.

If $s(m)>N$, then some odd prime power $p^e\parallel m$ satisfies
$e>c_p(N)\ge2N$, so $p^{2N+1}\mid m$. Consequently, for any positive
integer $M$,
\[
\frac{\#\{m\le M:s(m)>N\}}{M}
\le\sum_{p\text{ odd}}p^{-(2N+1)}.
\]
The sum of these bounds over $N\ge1$ converges. Dominated convergence
for the counting averages now proves \eqref{eq:mean}.
The enclosure follows from the rational computations and tail estimates
below, whose numerical method is justified in Appendix~\ref{app:certificates}.
\end{proof}

Let
\[
\mathcal M_K=1+\sum_{N=1}^{K}(1-D_N).
\]
The proof of Theorem~\ref{thm:tail} gives a particularly simple certified tail:
\begin{equation}\label{eq:meantail}
\frac98\,3^{-(2K+3)}
\le\mathcal M-\mathcal M_K
\le\frac98\left((1-2^{-(2K+3)})\zeta(2K+3)-1\right).
\end{equation}
For the lower bound, sum the mandatory prime-$3$ contributions. For the
upper bound, sum the union bound over primes and use
$(1-p^{-2})^{-1}\le9/8$ for every odd prime. Enlarging the sum from odd
primes to all odd integers at least $3$ gives exactly the displayed
zeta expression.

Taking $K=10$ and evaluating the finitely many zeta factors with rational
interval arithmetic yields
\[
\begin{split}
1.054025810088371518561523366670
&<\mathcal M,\\
\mathcal M&<
1.054025810088371612974596863433.
\end{split}
\]
The slightly wider decimal bounds in \eqref{eq:meaninterval} are convenient
to quote. They certify the rounded value
$\mathcal M\approx1.05402581008837$ without relying on floating-point
rounding or an unproved estimate of a zeta-series remainder.

\section{Computational verification and research scope}
\label{sec:verification}

\subsection{What was actually run}

The accompanying Python programs use only the standard library. The
verification run used Python 3.13.5; the source is written for Python 3.9
or later. The exact range checks recorded in
\texttt{data/verification.json} are summarized in
Table~\ref{tab:verification}. Ordinary integer and rational arithmetic is used
throughout; the density enclosures do not use floating-point values.

\begin{table}[htbp]
\centering
\caption{Completed checks. Formula-derived data beyond the independently
checked range are distinguished explicitly.}
\label{tab:verification}
\begin{tabular}{@{}p{.73\textwidth}r@{}}
\toprule
Check & Count or range\\
\midrule
Exact secant numbers generated by the triangle & $0\le n\le2100$\\
Agreement with the independent binomial recurrence & $0\le n\le100$\\
Finite-moment evaluations against exact coefficients & 1900\\
Moduli checked for period and onset consistency & $1\le m\le1000$\\
Equalities at the proposed period shift & 130793\\
Proper-divisor period candidates rejected by witnesses & 2432\\
All-index Euler-factor defect checks & 651\\
Huge-index reductions, including indices above $10^{100}$ & 18\\
Arithmetic checks of the two lifting-mode coefficients & 15\\
Moduli in the formula-generated CSV table & $1\le m\le10000$\\
Densities supplied with exact rational interval certificates & $1\le N\le10$\\
\bottomrule
\end{tabular}
\end{table}

For each modulus through $1000$, the code checks a complete proposed
period block and its shift, verifies a mismatch at $s(m)-1$ when the
preperiod is nonzero, and rejects each candidate $d(m)/\ell$ with $\ell$
a prime divisor of $d(m)$. These witnesses are stored explicitly.
Why are those candidates enough? Any proper divisor of $d(m)$ divides
at least one such $d(m)/\ell$, and a sequence with the smaller period
would also have the larger multiple as a period.

These finite checks support the implementation and catch mistakes, but
do not prove that a pattern observed on a finite prefix persists forever.
The universal statements rest on the mathematical proofs in the report.
No proof-assistant formalization was executed. The CSV entries for
$1001\le m\le10000$ are outputs of the proved formulas, not a claim of
independent coefficient-by-coefficient testing for that extra range.

The separate density program uses 96 terms of a positive Euler-transformed
eta series for each required zeta value. Its JSON file records exact
rational endpoints, outward-rounded decimal endpoints, the exceptional
prime factors, and the finite cutoff calculations. A reader can rerun
every reported finite calculation without network access or a computer
algebra system.

\subsection{What is settled, and what remains a different problem}

The OEIS assertion \eqref{eq:conjecture} is false, and its correct domain
of validity is exactly characterized. There is no remaining conjectural
step in the stated local periods, finite onset formula, global assembly,
density formulas, or mean formula. The proofs are conventional arguments
and are open to independent checking; the computations are supplementary.

A different question is whether the valuation tests in
\eqref{eq:oddpreperiod} admit a dramatically simpler closed form, or how
large the gap $\ceil{e/2}-s(p^e)$ can become. The complete effective
criterion does not by itself answer those structural questions.
Likewise, the remaining conjectures about the \emph{full} up/down
sequence in \cite{gulec} are not resolved here. Confusing that sequence
with its even subsequence would incorrectly turn a specialized result
into a claim about a different open problem.

The useful methodological point is narrow and concrete. A modular
sequence represented by finitely many power functions naturally divides
into unit modes and nilpotent modes. Pure periodicity can fail precisely
because the latter survive at a few initial indices. Identifying that
transient exactly converts an isolated counterexample into a complete
classification, and then into an arithmetic distribution problem.

\clearpage
\appendix
\section{A standalone counterexample checker}
\label{app:checker}

The following complete program requires Python 3.8 or later. It checks
only the finite disproof and does not depend on any theorem in the report.
The full archive uses Python 3.9 or later because it also uses
\texttt{math.lcm}.

\begin{lstlisting}[language=Python]
from math import comb

# sec(x) * cos(x) = 1, with integer coefficient arithmetic.
a = [1]
for n in range(1, 20):
    value = sum(
        (1 if j % 2 else -1) * comb(2*n, 2*j) * a[n-j]
        for j in range(1, n+1)
    )
    a.append(value)

print("a(19) =", a[19])
print("a(1) mod 27 =", a[1] % 27)
print("a(19) mod 27 =", a[19] % 27)

# phi(27) = 18. A period dividing 18 forces these equal.
if a[1] % 27 == a[19] % 27:
    raise RuntimeError("The claimed counterexample failed")
print("Conjecture disproved: 1 != 10 modulo 27.")
\end{lstlisting}

Output:
\begin{lstlisting}
a(19) = 23489580527043108252017828576198947741
a(1) mod 27 = 1
a(19) mod 27 = 10
Conjecture disproved: 1 != 10 modulo 27.
\end{lstlisting}

The same recurrence may be reduced modulo $27$ at every step, avoiding
large integers entirely. We retain the exact value here so that the
counterexample can be checked in either form.

\clearpage
\section{Rational certificates for the numerical constants}
\label{app:certificates}

For an integer $r>1$, let
\[
\eta(r)=\sum_{j\ge0}\frac{(-1)^j}{(j+1)^r}
       =(1-2^{1-r})\zeta(r).
\]
Define positive finite differences
\[
\Delta_{k,r}=\sum_{j=0}^k(-1)^j\binom{k}{j}(j+1)^{-r}.
\]
The elementary integral identity
\begin{equation}\label{eq:positiveintegral}
\Delta_{k,r}
=\frac1{(r-1)!}\int_0^1(-\log t)^{r-1}(1-t)^k\,dt
\end{equation}
shows $0<\Delta_{k,r}\le1$. Summing the geometric series in the integral
also proves the Euler transform
\begin{equation}\label{eq:eulertransform}
\eta(r)=\sum_{k\ge0}\frac{\Delta_{k,r}}{2^{k+1}}.
\end{equation}
Indeed, the sum of $(1-t)^k/2^{k+1}$ is $1/(1+t)$; integrating its
geometric expansion in $t$ gives the eta series. Positivity or dominated
convergence justifies the exchanges.

If
\[
H_{K,r}=\sum_{k=0}^{K-1}\frac{\Delta_{k,r}}{2^{k+1}},
\]
then both endpoints of
\begin{equation}\label{eq:zetacertificate}
\frac{H_{K,r}}{1-2^{1-r}}
\le\zeta(r)\le
\frac{H_{K,r}+2^{-K}}{1-2^{1-r}}
\end{equation}
are rational, and the enclosure follows because the omitted positive
terms sum to at most $2^{-K}$. The program forms the finite differences
recursively with \texttt{fractions.Fraction}, not by cancellation-prone
floating-point subtraction.

For $D_N$, the finite correction factor in \eqref{eq:densityzeta} is an
exact rational number. The reciprocal of the upper zeta endpoint gives
a lower density endpoint, and the reciprocal of the lower zeta endpoint
gives an upper density endpoint. Summing those bounds with the direction
reversed for $1-D_N$, then applying \eqref{eq:meantail}, gives the mean
certificate. Decimal conversions are rounded outwards using the Python
\texttt{decimal} module.

The certificate file includes unusually long integers because it stores
exact common denominators rather than a rounded approximation. On Python
versions with a decimal-digit limit for integer serialization, the
certificate program disables that limit only for these locally generated
integers. It does not parse untrusted large integer strings or fetch
external input.

\clearpage
\section{Archive contents and reproduction}
\label{app:archive}

The principal files are \texttt{article.tex}, \texttt{article.pdf},
\texttt{secant\_periodicity.py}, \texttt{verify.py}, and
\texttt{certify\_densities.py}. The \texttt{data/} directory contains
all tables, witnesses, logs, and rational certificates. A separate
\texttt{counterexample.py} reproduces Appendix~\ref{app:checker}.
The file \texttt{oeis\_correction.txt} is a proposed correction note for
human review; nothing has been submitted to the OEIS.

From the archive directory, run
\begin{lstlisting}[language=bash]
python counterexample.py
python verify.py
python certify_densities.py
latexmk -pdf -interaction=nonstopmode -halt-on-error article.tex
\end{lstlisting}
The first three commands need no third-party Python packages. Building
the PDF requires a LaTeX installation with the packages listed in the
source, including New PX text and math fonts. Font files are not included
in the archive. A \texttt{Makefile}, a PowerShell build script, and a
source-provenance note are included for convenience.

\begin{thebibliography}{9}
\addcontentsline{toc}{section}{References}
\bibitem{oeis}
The OEIS Foundation Inc., \emph{The On-Line Encyclopedia of Integer
Sequences}, entry A000364, Euler or secant numbers. In particular,
Peter Bala's conjecture comment dated 8 May 2023 and the period comment
attributed to Allen Stenger dated 3 August 2020. Consulted
19 September 2026. \url{https://oeis.org/A000364}.

\bibitem{knuth}
D.~E. Knuth and T.~J. Buckholtz,
Computation of tangent, Euler, and Bernoulli numbers,
\emph{Mathematics of Computation} \textbf{21}, no.~100 (1967), 663--688.
\url{https://doi.org/10.1090/S0025-5718-1967-0221735-9}.

\bibitem{metsankyla}
T. Mets\"ankyl\"a,
Note on Kummer's congruences for Euler numbers,
\emph{The Fibonacci Quarterly} \textbf{52}, no.~2 (2014), 160--162.
\url{https://www.fq.math.ca/Papers1/52-2/Metsankyla.pdf}.

\bibitem{sun}
Z.-W. Sun,
On Euler numbers modulo powers of two,
\emph{Journal of Number Theory} \textbf{115}, no.~2 (2005), 371--380.
Preprint arXiv:math/0410085v3, 7 December 2004.
\url{https://arxiv.org/abs/math/0410085}.
\url{https://doi.org/10.1016/j.jnt.2005.01.001}.

\bibitem{ramassamy}
S. Ramassamy,
Modular periodicity of the Euler numbers and a sequence by Arnold,
\emph{Arnold Mathematical Journal} \textbf{3}, no.~4 (2017), 519--524.
\url{https://arxiv.org/abs/1712.08666}.

\bibitem{gulec}
B. G\"ule\c{c},
Modular periodicity of the Euler up/down numbers at odd prime powers,
arXiv:2608.27058v2, 3 September 2026.
\url{https://arxiv.org/abs/2608.27058}.
\end{thebibliography}

\end{document}
